% Outline: Evaluation (3 page cap, >=6 figures; narrative order = story acts)
% Numbers of record: SYNC-CHECKLIST-2026-07-26.md; claims->evidence map in outline.md.
% Figure inventory for the >=6 rule, at exactly six: fig:workload, fig:tau,
% fig:overhead, fig:gate, fig:regime, fig:block1. The cold-start plate was
% dropped to reach the 3-page budget; its numbers are all in the prose and
% Fig. 2 in Background carries the per-stratum view.
\section{Evaluation}
\label{sec:eval}

\subsection{E1: What Agent Requests Are Made Of}
Fig.~\ref{fig:workload} characterizes the capture. It covers 75 requests
across 25 tasks of a production coding agent.
Each tool in the native sets contributes about 2\,kB of schema
(1.8--2.1\,kB), giving 21\,kB for ten native tools; the denser 170-tool
configuration totals 147\,kB, resent whole every turn. That fixed per-turn
cost is an unstable share of the request: 17\% to 57\% within the native
configuration alone (median 23.5\%), and 67.0\% with 170 tools. Within a
call kind the tool set is byte-identical across turns, and the 75 requests
carry only three distinct sets. Against the ToolACE vocabulary, every
tool-bearing live request classified as S4.
\ifrubric\else
\begin{figure*}[t]\centering
  % Source: probes/agent-traces-2026-07-26/MANIFEST.md + schema-variability probe.
  \includegraphics[width=\linewidth]{workload.pdf}
  \caption{Captured agent traffic: across the captured tool sets schemas are
  18--66\% of request bytes, a
  third of requests carry no tools, and completion medians split 6 versus 58
  tokens, so the pooled median of 42 describes neither kind.}
  \label{fig:workload}
\end{figure*}
\fi

\subsection{E2: Schema Text vs.\ Scalars vs.\ Identity}
Fig.~\ref{fig:tau} gives the headline offline result. Reading the request as
text beats the scalar baseline family by $+0.19$~$\tau$ against the tuned
baseline of record (0.6302 vs.\ 0.4395; fixed form 0.4268). The prompt-only
control reaches 0.5865, so roughly 77\% of that gap comes from the text
encoder and 23\% from the schema text. Isolated, the schema is worth
$\Delta\tau_b = +0.0437$ (95\% session-clustered paired CI
$[+0.024, +0.066]$). Schema identity is nearly worthless: a
fingerprint-plus-lookup baseline gains $+0.008$, inside its own $\pm 0.035$
interval, because only 4.5\% of test requests share a fingerprint with
training.

\begin{figure}[tb]\centering
  \includegraphics[width=\linewidth]{ranking.pdf}
  \caption{Ranking quality by input family: prompt-plus-schema text exceeds
  the tuned scalar baseline by $\Delta\tau_b{=}+0.191$. Dots are the mean
  over 3 seeds; capped bars the 95\% bootstrap CI on seed 17 alone; the lane
  below each row the min--max across seeds, a tick where that spread is
  zero.}
  \label{fig:tau}
\end{figure}

The same models are evaluated on the two strata Fig.~\ref{fig:strata} clears
for use: S3 (partially new tools, $n{=}543$) and S4 (all-new tools,
$n{=}333$). Both
pre-registered criteria pass on both strata against the tuned baseline, and
schema-text $\tau$ differs by only 0.0075 between them (0.647 and 0.639,
against 0.630 overall). We ran no equivalence test, so this is consistent
with no degradation rather than a demonstration of it. What does not hold
flat is the \emph{margin}: schema text leads the tuned baseline by $+0.25$
on S3 but only $+0.14$ on S4, because the scalar baseline itself scores
higher where no tool is familiar ($\tau_b$ 0.50 against 0.40). Reading the
schema keeps its own quality on wholly unseen tools; it does not widen its
lead there.

\subsection{E3: Deployment Cost and Two Failed Rescues}
Fig.~\ref{fig:overhead} prices the admission path the scorer sits on.
The Ranker costs 722\,ms at p99 on production CPU settings at the deployed
concurrency of eight, 48$\times$ the gateway's 15\,ms contract and still
8$\times$ over it served serially. Two pre-registered rescues fail.
\emph{Un-truncation}: 71\% of schemas exceed the 512-token cap, yet a
two-tower encoder reading the whole schema moves $\tau$ by $+0.0008$ (S3)
and $-0.0112$ (S4), intervals overlapping. \emph{Caching}: precomputing the
schema tower yields $1.9\times$, still 25$\times$ over contract at
concurrency 8 (4.3$\times$ serial). On the serving GPU with fp16
micro-batching, scoring costs 37.9\,ms at p50, $6.0\times$ faster than
unbatched GPU inference and
$16.6\times$ faster than the CPU path at the same percentile (628.8\,ms
p50). Enforcing a budget end to end is a different measurement, and we
report it over every request the gateway handled rather than over the scored
subset: under the documented 15\,ms default, 801 of 2{,}315 requests
(34.6\%) fail open into arrival order, silently; at 50\,ms that falls to
2.0\% and at 75\,ms to zero. The denominator includes requests the gate
never scores at all --- the zero-tool third of the traffic and the strata it
abstains on --- so 34.6\% is the share of \emph{traffic} left unordered, not
the share of scoring attempts that miss the budget, which is necessarily
larger.
\begin{figure*}[tb]\centering
  \includegraphics[width=\linewidth]{overhead.pdf}
  \caption{Admission-time scoring versus direct: $+$737\,ms TTFT and
  $+$840\,ms TTLT over the whole added path. Panel (a) is 150 paired
  requests and (b) the 118 of them with equal output tokens; violin bodies
  and $y$-axes are trimmed to p5--p95 of each arm, and intervals are
  2{,}000$\times$ bootstrap percentile CIs.}
  \label{fig:overhead}
\end{figure*}

\subsection{E4: Assigned versus Realized Confidence}
Fig.~\ref{fig:gate} contrasts the confidence each rule assigns with the
ranking quality it actually delivered.
The placeholder constant confidence (0.9) overstates realized ranking
quality on every stratum, by $+0.25$ to $+0.46$; the lower CI bound of
validation $\tau$ per stratum, abstaining where it cannot be measured, did
not overstate on any stratum here: its closest approach to overstating was
$-0.016$ on S4, and its largest understatement $-0.068$ on S3. S2 is where
the Ranker fails: familiar tools in new combinations realize $\tau = 0.44$,
against 0.60 on S1, 0.65 on S3 and 0.64 on S4 --- 0.16 to 0.21 below every
other stratum. At $n{=}78$ that sits below our reporting threshold, so we
read it as a design signal rather than a result, which is why
Fig.~\ref{fig:strata} plots neither it nor S1's.
\begin{figure*}[tb]\centering
  \includegraphics[width=\linewidth]{gate.pdf}
  \caption{Assigned versus realized confidence: the placeholder constant
  overstates on every stratum; Rule~C did not on this split, at the cost of
  abstaining where it cannot measure.}
  \label{fig:gate}
\end{figure*}

\subsection{E5: Serving-Level Scheduling}
We replayed 2{,}075 requests open-loop at 3.96 requests/s, 90\% of the
highest rate the stock scheduler sustained under a 95\% completion criterion
(4.4\,requests/s). Six schedulers share one arrival schedule: stock
first-come-first-served, a shim wrapping it unchanged, and four on our own
base, \emph{PolicyFCFS} (arrival order, no scorer), \emph{PromptLengthSJF}
(prompt length), \emph{PureLTR} (Ranker score) and \emph{GatedRuleC} (Ranker
score where the gate vouches). The pre-registered endpoint is pooled mean
time-to-last-token (TTLT), compared as paired-mean ratios with hierarchical
bootstrap intervals against a pre-declared 3\% margin; time-to-first-token
(TTFT) is reported only where the admission path is priced.

The attribution control comes first (Fig.~\ref{fig:block1}).
\emph{PolicyFCFS} and stock serve the same requests in the same arrival
order and differ only in scheduler implementation; PolicyFCFS records a
paired mean-TTLT ratio of $0.560$ ($[0.540, 0.581]$), so that 44\% reduction
arrives with the switch of implementation itself, none of it contributed by
ordering. Interposition does not explain it: the shim lands within 2\,ms of
stock (4903 vs.\ 4905\,ms). What does differ is the base class: the policy
arms subclass vLLM's \texttt{AsyncScheduler}, which overlaps scheduling with
model execution, while stock and the shim run the synchronous one --- the
likeliest mechanism, though the one-line ablation confirming it did not make
our budget.

Against the same-base control, ordering does not move the endpoint here. The
superiority test fails: GatedRuleC versus PromptLengthSJF gives $1.008$
($[0.979, 1.041]$), an interval containing one. Safety passes: against
PolicyFCFS the ratio is $0.996$ ($[0.970, 1.023]$), inside the 3\% margin,
and with an empty queue the gate itself costs at most 1.4\% (GatedRuleC against
PureLTR, $1.003$, $[0.991, 1.014]$; Sec.~V-F).

The order logs quantify how often ordering could act and show that the
policies acted when able. Across 77{,}000 scheduling steps per arm the
waiting queue was empty at p90 and held two waiting requests in under 1\% of
steps; 11.9--12.5\% of requests first entered a queue of depth two or more.
Expressing an order needs two requests that were \emph{already} waiting to
meet in the same step. That happened in 293--324 steps per arm, and the
policies rewrote 262--307 of them (87--95\%), while \emph{PolicyFCFS}
rewrote none. Ordering was executed on nearly every opportunity it had and
mean TTLT did not move. This round did not isolate why: three mechanisms are
consistent with it, the replay's compressed job-size variance, a two-deep
queue where a swap moves a request by at most one slot, and residual ranking
error at $\tau{=}0.63$. Round~A separates none of them; E6 does.

A second round replays all six arms with vLLM prefix caching enabled and
nothing else changed. The stock arms bridge the sessions, their on/off
ratios agreeing to four decimal places (0.8291 vs.\ 0.8290), and the
engine's log confirms the cache took effect (62.1\% peak hit rate). Caching
cuts pooled mean TTLT by 17.1\% on the stock pair and 22--23\% on the four
policy-base arms. The policy intervals overlap the bridge's, so the larger
gain is suggested rather than established, and the experiment did not
isolate the schema's share of it.

A third round at 6.4\,requests/s in three blocks of cold engine launches
fails the same check (exposure only 14.3--15.0\%, \emph{not one} carry-over
step across twelve launches), so it is non-diagnostic for superiority. On
those independent launches GatedRuleC gives $1.001$ ($[0.988, 1.015]$),
$1.006$ ($[0.993, 1.020]$) and $1.006$ ($[0.995, 1.018]$) against
PromptLengthSJF, PolicyFCFS and PureLTR, all inside the 3\% margin.

\subsection{E6: Ordering Under Contention}
None of those rounds formed a queue: the running-batch cap sat at vLLM's
default of 256 slots while only eleven to nineteen requests were in flight.
E6 caps it at 16 and repeats the question at 6.4 requests/s with prefix
caching on and one cold launch per arm. Three conditions therefore differ
from round~A besides the cap --- rate, caching and launch structure --- so
the comparison across panels is between regimes rather than between two
otherwise identical runs; what is held exactly is the workload, the arrival
offsets, the checkpoint, and the four policies being contrasted, which is
what makes the ordering comparison \emph{within} each regime clean. The cap
is an instrument, not a recommendation. A scheduler can only act where
requests wait, and what we could vary at our budget was the number of slots
rather than the arrival rate needed to fill 256 of them; a deployment
reaches the same condition whenever offered concurrency exceeds the slots
it has, whether because the cap is low or because long contexts exhaust the
KV cache first. What E6 tests is therefore the regime, not the number: the
claim is that ordering pays once a queue exists, and the cap is how we made
one. The
manipulation check passes this time: 57.5--63.6\% of requests entered a
queue of depth two or more, p90 depth reached 8 to 14, and 5{,}914 to
7{,}121 steps per arm carried two requests that had both waited earlier,
against 293--324 in round~A. The policies rewrote 1{,}817 to 6{,}607 of
them, and \emph{PolicyFCFS} again rewrote none.

Ordering now moves the endpoint (Fig.~\ref{fig:regime}), and the
pre-registered primary test passes for the first time: GatedRuleC versus PromptLengthSJF gives $0.966$
($[0.935, 0.998]$), an upper bound below one. What pays is the learned
signal. PureLTR beats arrival order by 4.2\% ($0.958$, $[0.944, 0.972]$),
while the free length heuristic runs 8.5\% \emph{worse} than doing nothing
($1.085$, $[1.051, 1.122]$), so rendered prompt length is a poor proxy for
output length; over 2{,}025 shared sessions pooled mean TTLT is 2.797\,s for
PureLTR against 3.169\,s for PromptLengthSJF. The two regimes also separate
the gate's cost:
GatedRuleC against PureLTR was $1.003$ ($[0.991, 1.014]$) with an empty
queue and is $1.095$ ($[1.079, 1.112]$) here, and against PolicyFCFS it is
$1.049$ ($[1.033, 1.065]$), outside the 3\% margin. Abstention pins S1, S2
and zero-tool requests in their arrival slots, and once the queue is deep
those pinned requests block the ones behind them. Abstaining is free when
nothing is waiting and costly when something is.

\begin{figure*}[tb]\centering
  % Source: runs/block1-main/matrix (default cap) and
  % runs/capped-batch/matrix (cap 16); same workload, same arrival offsets,
  % same checkpoint. Ratios and intervals are the paired bootstrap.
  \includegraphics[width=\linewidth]{regime.pdf}
  \caption{One engine setting decides the paper's answer: the same four
  policies on the same workload, with the default running-batch cap no pair
  separates, and with the cap reduced every pair does. Both panels replay the
  same 2{,}025 sessions and each arm resamples within one engine launch (3
  replays in (a), 1 in (b)), so neither panel's intervals carry machine-level
  replication.}
  \label{fig:regime}
\end{figure*}
\begin{figure*}[tb]\centering
  % Source: runs/block1-main/matrix, six arms x 3 replay repetitions
  % (round A only; the generator reads matrix/, not matrix-round-b/),
  % 2075-request replay.
  \includegraphics[width=\linewidth]{block1.pdf}
  \caption{Serving-level scheduling at the default batch cap: the
  $1.78\times$ gap between stock and \emph{PolicyFCFS} comes from the
  scheduler implementation, not from ordering, and the policies rewrote
  87--95\% of the steps that could express an order. The paired ordering
  ratios for this regime, and the 3\% margin they are judged against, are in
  Fig.~\ref{fig:regime}(a).}
  \label{fig:block1}
\end{figure*}
